\documentclass[10pt,letterpaper]{article}
\usepackage{cornell-notes}

\title{Clause 33.03.03.04: Non Member Functions}
\author{}
\date{}
\setDocTitle{Clause 33.03.03.04: Non Member Functions}
\setDocSubtitle{C++ 2024}
\setDocOwner{C++ 2024 Cornell Notes}
\setDocDate{\today}
\setCornellCollection{C++ 2024 Cornell Notes}
\setCornellUnitType{Clause}
\setCornellUnitNumber{33.03.03.04}
\setCornellUnitTitle{Non Member Functions}
\hypersetup{pdftitle={Clause 33.03.03.04: Non Member Functions},pdfsubject={Cornell notes},pdfauthor={}}

\begin{document}
\maketitle
\begin{CornellOverviewBox}{Overview}
\textbf{Classification:} Standard library - Normative\\[2pt]
\textbf{Learning objective:} Understand the rules, terminology, and semantic consequences associated with Non-member functions.
\end{CornellOverviewBox}\begin{CornellSummaryBox}{Summary}
{\sffamily\bfseries\color{Accent} Summary}\par\vspace{3pt}Section 33.3.3.4 focuses on Non-member functions. Apply the specified interface together with its mandates, effects, result conditions, exception behavior, and complexity constraints. When applying it, preserve the distinction between mandatory requirements, recommendations, permissions, and behavior deliberately left undefined, unspecified, or implementation-defined.
\end{CornellSummaryBox}

\end{document}
